\begin{table}[H]
\caption{Среднее время отклика на основные действия в системе анализа экзопланет\label{time}}
\centering
\begin{tabular}{|p{9cm}|r|}
    \hline
    \textbf{Функциональность} & \textbf{Среднее время отклика, сек} \\
    \hline
    Валидация параметров экзопланеты & 0.45 \\
    \hline
    Расчет индекса подобия Земле (ESI) & 0.82 \\
    \hline
    Построение 3D-модели экзопланеты & 1.35 \\
    \hline
    Генерация графиков и диаграмм & 1.12 \\
    \hline
    Поиск похожих экзопланет в базе & 0.98 \\
    \hline
    Расчет вероятности наличия воды & 0.89 \\
    \hline
    Анализ атмосферных параметров & 0.93 \\
    \hline
    \textbf{Среднее значение по всем операциям} & \textbf{0.93} \\
    \hline
\end{tabular}
\end{table}

Исходя из результатов тестирования в таблице \ref{time}, можно сделать вывод, что время 
отклика на все виды запросов находится в пределах нормы.

\textbf{Оценка ресурсных показателей программного средства}

Для оценки ресурсных показателей системы анализа экзопланет использовались встроенные инструменты браузера Google Chrome — в частности, Chrome DevTools Performance Profiler. Он позволяет в реальном времени отслеживать использование ресурсов при взаимодействии пользователя с интерфейсом: от ввода параметров до финального отображения результатов анализа.

Оптимизация производительности клиентской части системы предоставляет ряд важных преимуществ, включая снижение времени расчетов, повышение отзывчивости интерфейса и комфортность использования приложения на различных устройствах.

Performance Profiler в Chrome позволяет:
\begin{itemize}
    \item визуализировать процесс расчета параметров обитаемости;
    \item определить этапы построения 3D-моделей и графиков;
    \item отслеживать время отклика на изменение параметров;
    \item измерить метрики, характеризующие восприятие производительности.
\end{itemize}

Информация, собираемая при профилировании:
\begin{itemize}
    \item LCP (Largest Contentful Paint) — фиксирует время отображения 3D-модели и основных графиков;
    \item CLS (Cumulative Layout Shift) — показывает стабильность интерфейса при обновлении данных;
    \item INP (Interaction to Next Paint) — измеряет задержку при интерактивном взаимодействии с 3D-моделью.
\end{itemize}

После проведения анализа производительности интерфейса системы были получены следующие показатели:

\begin{itemize}
    \item LCP: 1.35 секунды — время построения 3D-модели экзопланеты
    \item CLS: 0.02 — минимальные смещения при обновлении графиков
    \item INP: 75 мс — быстрый отклик на вращение 3D-модели
\end{itemize}

\textbf{Оценка показателей надёжности программного средства}

Для проверки надёжности системы анализа экзопланет была применена упрощённая оценка надёжности, основанная на количественной фиксации сбоев при многократном использовании системы.

Для оценки надёжности программного средства использовалась следующая формула:

\begin{equation}
P = 1 - \frac{n}{N}, \label{eq:reliability}
\end{equation}

где $n$ — количество зафиксированных сбоев,  
$N$ — общее число запусков/испытаний системы.

В ходе функционального тестирования было выполнено 100 последовательных запусков ключевых сценариев:
\begin{itemize}
    \item ввод параметров экзопланеты
    \item расчет индекса подобия Земле
    \item построение 3D-модели
    \item генерация графиков
    \item поиск похожих экзопланет
\end{itemize}

В процессе тестирования было зафиксировано 3 нестабильных поведения:
\begin{itemize}
    \item один случай некорректного построения 3D-модели
    \item один случай задержки при расчете индекса ESI
    \item один случай ошибки при поиске похожих экзопланет
\end{itemize}

Тогда, согласно формуле:

\begin{equation}
P = 1 - \frac{3}{100} = 0{,}97. \label{eq:reliability-result}
\end{equation}

Полученное значение $0{,}97$ свидетельствует о высокой степени надёжности системы анализа экзопланет. Это означает, что система корректно функционирует в 97\% случаев при стандартных пользовательских сценариях.

Основные аспекты надёжности системы:
\begin{itemize}
    \item точность расчета параметров обитаемости
    \item стабильность визуализации данных
    \item корректность поиска похожих экзопланет
    \item устойчивость работы 3D-рендеринга
\end{itemize}

\newpage
\subsection{Вывод}

В ходе разработки интеллектуальной системы анализа обитаемости экзопланет были использованы современные веб-технологии, включая Three.js для 3D-визуализации и Plotly для построения графиков. Были проведены работы по разработке математических моделей, алгоритмов анализа и визуального представления результатов. Основной целью проекта было создание инструмента, позволяющего исследователям эффективно оценивать потенциальную обитаемость экзопланет на основе их физических параметров.

В процессе разработки системы были учтены требования к точности расчетов, надежности и удобству использования. Были реализованы основные функциональные возможности, включая:
\begin{itemize}
    \item валидацию входных параметров экзопланет
    \item расчет индекса подобия Земле (ESI)
    \item анализ вероятности наличия жидкой воды
    \item построение интерактивных 3D-моделей
    \item генерацию наглядных графиков и диаграмм
    \item поиск и сравнение похожих экзопланет
\end{itemize}

После завершения разработки система была подвергнута тщательному тестированию, в ходе которого проверялась точность расчетов, корректность визуализации и отсутствие ошибок. Благодаря тестированию были выявлены и устранены возможные недочеты, что повысило надежность и качество системы.

В результате успешной реализации была достигнута основная цель проекта. Разработанная система предоставляет исследователям эффективный инструмент для анализа потенциальной обитаемости экзопланет. Она может быть использована в научных исследованиях и образовательных целях, что делает ее полезной для широкого круга специалистов в области астрономии и планетологии.

В целом, выполнение данного пункта дипломного проекта подтверждает успешное завершение разработки системы анализа обитаемости экзопланет и ее готовность к практическому применению. Разработанная система представляет собой значимый вклад в развитие инструментов для исследования экзопланет и поддерживает дальнейшее развитие методов оценки их потенциальной обитаемости. 